\section*{ВВЕДЕНИЕ}
\addcontentsline{toc}{section}{ВВЕДЕНИЕ}

\textbf{Общая информация о проекте:}\\
\textbf{Название проекта:} «ЭкоМедиа: Медиа-проектирование в экологии человека»

\textbf{Актуальность темы.}\\
В двадцать первом веке человечество столкнулось с глобальными вызовами, обусловленными стремительной урбанизацией, цифровизацией и масштабным производством синтетических материалов. Жизнь в современном мегаполисе неразрывно связана с постоянным воздействием множества неблагоприятных факторов: загрязнением атмосферного воздуха, ухудшением качества потребляемой питьевой воды, повсеместным распространением микропластика в пищевых цепях, а также нарушением биологических ритмов человека из-за избыточного использования цифровых устройств. Экология человека, как междисциплинарная наука, исследует эти закономерности и разрабатывает методы минимизации вредного воздействия окружающей среды на здоровье.

Одной из главных проблем современного общества является низкий уровень информированности населения, особенно молодежи, о скрытых экологических и цифровых угрозах. Большинство людей не задумываются о том, как их повседневные привычки — от выбора посуды для разогрева пищи до просмотра ленты новостей перед сном — медленно, но верно разрушают их здоровье. В связи с этим возникает острая необходимость в создании качественного, современного и доступного медиа-контента, который способен привлечь внимание к проблемам экологии человека и предложить простые, практичные пути их решения. Именно это обуславливает высокую актуальность проекта «ЭкоМедиа».

\textbf{Объект и предмет разработки.}\\
Объектом данного проекта является процесс экологического просвещения и формирования навыков цифровой гигиены у жителей современных городов. Предметом проекта выступает разработка комплексного медиа-ресурса (веб-сайта) и интерактивного программного продукта (Telegram-бота), способствующих популяризации здорового образа жизни.

\textbf{Цели и задачи проекта:}\\
\textbf{Цель проекта} заключается в повышении уровня экологической грамотности и формировании здоровых, осознанных привычек у целевой аудитории (преимущественно студентов и молодежи) в условиях городской среды через создание инновационного медиа-ресурса и интерактивных инструментов.

Для достижения поставленной цели в рамках проектной (учебной) практики были сформулированы и решены следующие \textbf{задачи}:
\begin{enumerate}
    \item Провести всесторонний анализ ключевых факторов риска для здоровья человека в городе, уделив особое внимание проблеме микропластика, качеству водопроводной воды и влиянию синего света от экранов гаджетов на физиологию сна.
    \item Осуществить поиск, систематизацию и проверку научной информации по выбранным направлениям экологии человека для подготовки достоверного контента.
    \item Спроектировать и разработать статический веб-сайт с современным адаптивным дизайном (с применением концепции glassmorphism) для презентации собранных материалов, статей и видеороликов.
    \item Настроить инфраструктуру коллективной разработки с использованием системы контроля версий Git и платформы GitHub.
    \item В рамках вариативной части задания спроектировать и разработать Telegram-бота на языке программирования Python, реализующего функции интерактивного тестирования знаний пользователей, оценки уровня их цифровой гигиены и выдачи ежедневных экологических советов.
    \item Задокументировать все этапы жизненного цикла проекта, включая архитектуру разработанного программного обеспечения, и составить итоговый технический отчет.
\end{enumerate}

Выполнение перечисленных задач направлено на создание полностью функционального продукта, готового к практическому применению и дальнейшему масштабированию в рамках просветительских кампаний.
